\documentclass[10pt,a4paper]{amsart}
\usepackage[margin=19mm]{geometry}
\usepackage{amsmath,amssymb,mathtools}
\usepackage[scr=rsfs,cal=euler]{mathalfa}
\usepackage{xfrac}
\usepackage[unicode,hidelinks,bookmarks=false]{hyperref}
\newcommand{\ie}{i.e.\ }
\newcommand{\cf}{cf.\ }
\newcommand{\resp}{respectively\ }
\newcommand{\Subsection}{\subsection}
\providecommand{\nobreakdash}{\nobreak\hbox{-}}
\setlength{\emergencystretch}{2em}
\multlinegap0pt
\usepackage{bm}
\usepackage{bbm}
\usepackage{dsfont}
\usepackage{xspace}
\usepackage{cancel}
\usepackage{mathtools}
\usepackage{amsthm}
\makeatletter
\@ifundefined{definition}{\newtheorem{definition}{Definition}}{}
\@ifundefined{lemma}{\newtheorem{lemma}[definition]{Lemma}}{}
\makeatother
\newcommand{\ds}{\,\mathrm{d}s}
\newcommand{\ite}{i.\,e.\xspace}
\title[An improved bound for the ground state of a Schr\"odinger op]{An improved bound for the ground state of a Schr\"odinger operator on a loop}
\author{Helmut Linde}
\date{Source: arXiv:2504.20229v2 (CC BY 4.0)}
\begin{document}
\maketitle

\noindent\emph{Source and licence.} An improved bound for the ground state of a Schr\"odinger operator on a loop. Version arXiv:2504.20229v2 (CC BY 4.0), distributed under the Creative Commons Attribution 4.0 International licence, as recorded in the arXiv licence metadata for this exact version and shown on the abstract page https://arxiv.org/abs/2504.20229v2. Full rights notice: licenses/arxiv-2504.20229-CC-BY-4.0.txt.\par\smallskip

\noindent\textit{Editorial scope.} Counted unit: the Lemma whose source label is \texttt{lem:facts\_about\_I} in the article (source0.tex lines 412--418, source label \texttt{lem:facts\_about\_I}), delivered together with the article's own complete original proof of it (source0.tex lines 420--438, one \texttt{proof} environment of 19 lines). No mathematics is written, repaired or re-derived: statement and proof are the article's own text line for line. Required context retained uncounted: EqLambda (lines 375--377); eq:def\_h (lines 380--382); eq:def\_I (lines 384--386). Every definition, display and label of those lines is retained unchanged. Omitted from this package and not redistributed: 1--374, 378--379, 383--383, 387--411, 419--419, 439--1833. Nothing inside a retained excerpt is omitted or altered: every retained body is delivered exactly as the article's lines are written, so no omission receipt of any kind accompanies this package. Citations: the retained excerpts carry no citation command at all, so no bibliography entry of the article is transcribed and no reference is redistributed. Reference adaptation: none was needed and none was made; every reference inside the delivered unit resolves inside it, so no unresolved reference or citation is produced. Printed slips kept literal: no printed word, symbol, hypothesis or number of the counted statement or of its proof was altered. Layout and class: the article's own class is \texttt{article} with packages including amsmath, amsthm, amssymb, biblatex, booktabs, cleveref, comment, eso-pic, graphicx, ifthen, inputenc, multicol, multirow, numprint; the wrapper is the lane's pinned \texttt{amsart} editorial wrapper at 10pt on A4, which restates the theorem-like environments the retained text needs and adds the article's own macro definitions that the retained text needs (\texttt{\textbackslash ds}, \texttt{\textbackslash ite}), with extra wrapper packages (bm, bbm, dsfont, xspace, cancel, mathtools). The delivered numbering is the unit's own. Assets: no retained span contains a figure environment, a graphics-inclusion command, a PDF-page inclusion, a table or a TikZ picture; no image file is copied into this package, the package declares no asset dependency and no image file is redistributed with it. Licence: version arXiv:2504.20229v2 of this article is distributed under the Creative Commons Attribution 4.0 International licence, as recorded in the arXiv licence metadata for this exact version and shown on the abstract page https://arxiv.org/abs/2504.20229v2. A full rights notice is delivered at licenses/arxiv-2504.20229-CC-BY-4.0.txt. Characters: every retained excerpt is pure ASCII, so the delivered engine, XeLaTeX with its default Latin Modern text font, reports no missing character.
\begin{equation}\label{EqLambda}
    E(x,y) = \frac{\int_\Omega \left(x'(s)^2 + y'(s)^2 \right)\ds}{\int_\Omega \left(x(s)^2 + y(s)^2 \right)\ds}.
\end{equation}

\begin{equation} \label{eq:def_h}
    h_t(s) :=  x(s) \sin t - y(s) \cos t
\end{equation} 

\begin{equation} \label{eq:def_I}
    I(t) := \frac{\int_\Omega h'_t(s)^2 \ds}{\int_\Omega h_t(s)^2 \ds},
\end{equation}%

\begin{lemma} \label{lem:facts_about_I}
    Either 
    
    $I(t)$ is constant and then $\lambda_\Gamma\ge 1$, or 
    
    $I(t)$ has exactly one pair of maxima at a distance of $\pi$ from each other and exactly one pair of minima also at a distance of $\pi$ from each other and it is strictly monotonous in between those extrema.
\end{lemma}

\begin{proof}
    If $I(t) = I$ is constant then $I(t) \ge 1$ for all $t\in\Omega$ since we already know that this inequality holds if $t$ is any of the at least six critical angles of $\Gamma$. In that case we also have
    \[I = I(0) = \frac{\int_\Omega y'(s)^2\ds}{\int_\Omega y(s)^2\ds} = I(\frac{\pi}{2}) = \frac{\int_\Omega x'(s)^2\ds}{\int_\Omega x(s)^2\ds}\]
    and therefore by (\ref{EqLambda}) we get  $\lambda_\Gamma = I \ge 1$.
    
    So assume now that $I(t)$ is not constant. Since $I(t)$ is $\pi$-periodic, the following analysis can be restricted to the case $t \in (-\frac\pi2,\frac\pi2)$ where $\cos t>0$ plus the edge case $t = \frac\pi2$.
    Put (\ref{eq:def_h}) into (\ref{eq:def_I}), simplify the fraction with $\cos^2 t$ and write shorthand $\tau(t) = \tan t$ to obtain
    \begin{equation*}
         I(t) = \frac{(\int_\Omega x'^2\ds) \tau^2 -(\int_\Omega 2x'y' \ds) \tau +(\int_\Omega y'^2 \ds)}{(\int_\Omega x^2\ds) \tau^2 -(\int_\Omega 2xy \ds) \tau +(\int_\Omega y^2 \ds)}.
    \end{equation*}
    By the quotient rule the derivative of $I(t)$ can be written as
    \[I'(t) = \frac{\tau' P_3(\tau)}{P_4(\tau)}\]
    where $P_3(\tau)$ is a polynomial in $\tau$ of up to third degree and $P_4(\tau)$ is a polynomial in $\tau$ of exactly fourth degree.

    Being $\pi$-periodic but not constant, $I(t)$ must have a positive even number of extrema on $(-\frac\pi 2,\frac\pi 2]$, alternating between maxima and minima.
    The case of two extrema is claimed in the lemma, so it only remains to show that the number of extrema cannot be four or higher.

    For $I'(t)$ to have at least four zeros on $(-\frac\pi 2,\frac\pi 2]$, the polynomial $P_3$ must be of third degree to have three zeros on $(-\frac\pi 2,\frac\pi 2)$ and there must be a fourth zero at $t = \frac\pi 2$. Yet this fourth zero cannot exist: Being of third degree, $P_3(\tau)$ behaves like some (non-zero) constant times $\tau^3$ as $\tau \rightarrow \infty$ (\ite $t \rightarrow \frac\pi 2$). It is multiplied by $\tau'$ which increases like $\tau^2$ as $\tau \rightarrow \infty$. Their product then behaves like $\tau^5$. Their quotient with $P_4$, \ite the derivative $I'(t)$, can thus not converge to zero as $\tau\rightarrow\infty$. The continuity of $I'(t)$ then ensures that $I'(\frac\pi 2) \neq 0$.
\end{proof}

\end{document}
